\section{Physical clearance caustics and positive source localization}
\label{sec:v62-clearance-caustics}
We now locate the remaining loss of transversality by a geometric set
in parameter--roof space. For each fixed collision count and a positive
incidence threshold, its fiber at a fixed radius is finite. Away from
this set a semialgebraic separation inequality gives a quantitative
coarea denominator. The exact positive source over a neighborhood of
the set remains in the raw-error formula.

\subsection{The physical-side compact graph}
Fix $m\ge2$ and $0<\chi<1$. In every true center word, recenter at a
possible clearance witness $y=T_R^{j+1}x$, as in
\eqref{eq:v62-fiber-objects}, and require
\begin{equation}\label{eq:v62-incidence-cap}
 c_i(x)\ge\chi\quad(0\le i\le m).
\end{equation}
Take the closure of this regular physical graph with its parameter
$R\in I$, rather than replacing its strict first-hit inequalities by
an arbitrary nonphysical word. Retain the compact one-flight envelope
\[
 R\le |Z_d(y)|\le R+s_0,\qquad L_d(y)\ge\ell_*>0,
\]
where $\ell_*$ is the constant in
Lemma~\ref{lem:v49-physical-envelope}. Let $\mathcal K_{m,\chi}$ be
the finite disjoint union of these closures over words, marks,
candidates and compact rational chart pieces. A fixed finite chart
cover can be chosen with interiors covering every collision angle.
The expression $\mathcal J=\det D(F,Z)$ below is computed in its
particular marked chart. Its zero set is independent of that choice.

The positive incidence threshold makes the chosen contact roots simple,
including on these closures. The implicit contact equations therefore
supply continuous jets there. These are limits of physical-side jets;
the construction does not continue a source trajectory across a
competing-hit seam. Uniform upper and lower flight lengths, bounded
contact coordinates and velocities make $\mathcal K_{m,\chi}$ compact.
It is semialgebraic by the auxiliary graph construction and closure.

Define the clearance caustic by
\begin{equation}\label{eq:v62-caustic-set}
 \mathcal C_{m,\chi}
   =\{(R,F(y)):\ y\in\mathcal K_{m,\chi},\ |Z(y)|=R,
                                        \ \mathcal J(y)=0\}.
\end{equation}
Different word copies are allowed to have the same image. In particular
all roof-critical seam contacts are retained in this definition.

\begin{theorem}[Finite physical caustic fibers]
\label{thm:v62-caustic-finite}
The set $\mathcal C_{m,\chi}$ is compact and semialgebraic. There are
$C>0$, $A>1$, independent of $m,\chi,R$, such that
\begin{equation}\label{eq:v62-caustic-fiber-count}
 \#\{t:(R,t)\in\mathcal C_{m,\chi}\}\le CA^m
          \qquad(R\in I).
\end{equation}
Thus it is a finite union of parameter-dependent semialgebraic arcs
and points and has no interval in a vertical fiber.
\end{theorem}
\begin{proof}
The compact image theorem and Tarski--Seidenberg give the first claim.
For fixed $R$, the one-flight formulas of
\eqref{eq:v49-backward-distance} give
$\partial_\varphi Z_d=L_d\ge\ell_*$ on the relevant envelope.
Since $c=\cos\varphi\ge\chi$, this is a nonzero derivative also in
the marked $(\alpha,p)$ charts. Hence $dZ\ne0$ there.
On $Z=R$ or $Z=-R$, the condition $\mathcal J=0$ says that $dF$
is proportional to $dZ$. Every connected semialgebraic component of
this set is joined by piecewise $C^1$ paths; this follows from cell
decomposition, or \cite[Sections 2.3.4 and 3.2]{CosteSAG}. Along each
piece of a path, $dZ=0$ and therefore $dF=0$. Continuity at the
finitely many joining points shows that $F$ is constant on the whole
component. Isolated points also contribute only one value.

It remains to bound the number of components before projection.
First jets are included exactly as in
Lemma~\ref{lem:v62-jet-complexity}; all implicit coefficients are
nonzero under \eqref{eq:v62-incidence-cap}. To avoid a complexity claim about eliminating a closure operation,
use the following larger graph only for counting values. Keep the
contact equations, positive incidence cap, root-selection signs and
jet equations; replace the physical first-hit inequalities by their
weak versions and adjoin the seam and rank equations. This graph has
$O(m+1)$ variables and fixed-degree sign conditions and contains the
capped physical-side closure. Its selected contact roots are still
simple, so its jets are the unique implicit jets of those equations.
On every connected component of this larger critical graph the same
$dZ\ne0$, $dF\wedge dZ=0$ argument makes $F$ constant. Thus a subset
of such a component cannot create additional roof values, even if the
subset is disconnected. The bound \eqref{eq:v44-sign-bound} gives
$CD^m$ components of the larger graph, uniformly in the coefficient
$\chi$ and the parameter $R$. No measure or orbit is continued on
this counting overcover. Summing the exponentially
many words and possible marks changes only the exponential constant.
Since each component contributes at most one roof value, this proves
\eqref{eq:v62-caustic-fiber-count}. Semialgebraic cell decomposition
of the image, together with the finite fibers, gives the last assertion.
\end{proof}

\subsection{A denominator off the parameter--roof caustic}
For $v\in\mathcal K_{m,\chi}$ put
\[
 d(v)=\min\{1,\operatorname{dist}((R,F(v)),\mathcal C_{m,\chi})\},
 \qquad f(v)=\bigl||Z(v)|-R\bigr|+|\mathcal J(v)|.
\]
If the caustic is empty, define $d=1$.
The distance is taken in the two-dimensional parameter--roof space,
not to the fiber at the same radius. This convention retains contacts
at a parameter transition.

\begin{lemma}[Physical-side semialgebraic separation]
\label{lem:v62-caustic-separation}
For every fixed $m,\chi$ there are $C_{m,\chi}\ge1$ and an integer
$N_{m,\chi}\ge1$ such that
\begin{equation}\label{eq:v62-caustic-separation}
 d(v)^{N_{m,\chi}}\le C_{m,\chi}f(v)
                          \qquad(v\in\mathcal K_{m,\chi}).
\end{equation}
No bound on the growth of these constants as $m\to\infty$ or
$\chi\to0$ is asserted.
\end{lemma}
\begin{proof}
The jets, $f$ and $d$ are continuous semialgebraic functions on the
compact finite union of graph pieces. Semialgebraicity of distance
follows from its defining infimum and Tarski--Seidenberg. If $f=0$,
its parameter--roof image belongs to \eqref{eq:v62-caustic-set}, so
$d=0$. The compact semialgebraic \L ojasiewicz inequality
\cite[Theorem 2.12]{CosteSAG} gives \eqref{eq:v62-caustic-separation}
on each piece; finite maxima of constants and exponents combine them,
since $0\le d\le1$. If the caustic is empty, $f$ has a positive
minimum on the compact graph, and the conclusion follows directly.
\end{proof}

At an original first-clearance point $f\le2\varepsilon+|\mathcal J|$.
Consequently, for $0<h\le1$,
\begin{equation}\label{eq:v62-caustic-rank}
 \varepsilon\le\frac{h^{N_{m,\chi}}}{4C_{m,\chi}},\quad d\ge h
 \quad\Longrightarrow\quad
 |\mathcal J|\ge\frac{h^{N_{m,\chi}}}{2C_{m,\chi}}.
\end{equation}
This is a bound on the actual determinant, uniform in the radius for
the fixed $m,\chi,h$.

Define $b^{\varepsilon,\mathrm{cau}(\chi,h),w}$ to be the original
first-clearance source restricted to all $c_i\ge\chi$ and $d<h$.
The distance uses the selected witness chart through the common
parameter--roof image. The source with some $c_i<\chi$ is kept as
a separate positive source, even if the first physical defect was
clearance and the small incidence occurs later.

\begin{theorem}[Positive caustic localization at exact labels]
\label{thm:v62-caustic-localization}
There are $C>0$, $A>1$ such that, under the threshold condition in
\eqref{eq:v62-caustic-rank}, $0<2\varepsilon<s_0$ and $|w|\le M$,
\begin{align}\label{eq:v62-caustic-localization}
 &\operatorname*{ess\,sup}_t\sum_{n=1}^m\sum_{k\in\Z^2}
 \left|b^{\varepsilon,\mathrm{clr},w}_{n,k,m,R}(t)
       -b^{\varepsilon,\mathrm{cau}(\chi,h),w}_{n,k,m,R}(t)\right|
 \notag\\
 &\hspace{15mm}\le CMA^m\left((m+1)\chi+
                C_{m,\chi}\varepsilon h^{-N_{m,\chi}}\right).
\end{align}
The caustic term is a restriction of the original positive source when
$w=1$. This is an identity-based localization of that source on all
roofs, not a deletion of a set of roof values.
\end{theorem}
\begin{proof}
Split the original clearance source into: some $c_i<\chi$; all
$c_i\ge\chi$ and $d<h$; and all $c_i\ge\chi$ and $d\ge h$.
On a regular physical word $0<c_i\le1$, so
\[
 \one_{\{\min_i c_i<\chi\}}
       \le \sum_{i=0}^{m}\chi c_i^{-1}
       \le(m+1)\chi\,\mathcal I_{m,R}.
\]
Theorem~\ref{thm:v61-inverse-incidence-density} bounds the height
of the first part by $CMA^m(m+1)\chi$. This comparison does not
change the first defect's type or erase its later history.
For the third part, \eqref{eq:v62-caustic-rank} gives a lower
bound for the marked-chart determinant. Invariance under recentering
and the bounded rational-chart density let us apply the same area
formula and physical multiplicity proof as
Theorem~\ref{thm:v61-transverse-clearance}, now in that chart.
The total length of the signed clearance bands, summed over marks,
is $O(\varepsilon)$, so its height is at most
$CMA^m C_{m,\chi}\varepsilon h^{-N_{m,\chi}}$.
Their sum is the left side for $w=1$. Domination and the disjoint
exact-label partition prove the inserted statement. The retained middle
part is precisely the caustic source in the theorem.
\end{proof}

\begin{corollary}[The raw-error formula with the caustic source]
\label{cor:v62-caustic-raw-error}
For $\varepsilon=\varepsilon(B)$ satisfying
\eqref{eq:v62-caustic-rank}, uniformly over $R,n,k$,
\begin{align}\label{eq:v62-caustic-raw-error}
 \|P-G-m^2b^{\varepsilon,\mathrm{cau}(\chi,h),1}\|_\infty
 \le e_{B,m}+Cm^2A^m\left(\varepsilon+(m+1)\chi+
             C_{m,\chi}\varepsilon h^{-N_{m,\chi}}\right).
\end{align}
\end{corollary}
\begin{proof}
Use Theorem~\ref{thm:v51-positive-error}, the original incidence
height gain and Theorem~\ref{thm:v62-caustic-localization}. The
canonical finite arithmetic reference $G$ is not modified.
\end{proof}

\subsection{Size, but not height, of the retained caustic source}
\begin{lemma}[Uniform vertical measure of a caustic neighborhood]
\label{lem:v62-caustic-tube}
For fixed $m,\chi$ there are $D_{m,\chi}<\infty$ and
$\beta_{m,\chi}>0$ such that
\begin{equation}\label{eq:v62-caustic-tube}
 \sup_{R\in I}|\{t:\operatorname{dist}((R,t),\mathcal C_{m,\chi})<h\}|
       \le D_{m,\chi}h^{\beta_{m,\chi}},\qquad 0<h\le1.
\end{equation}
For the empty caustic the left side is zero.
\end{lemma}
\begin{proof}
First the left side tends uniformly to zero as $h\downarrow0$.
Otherwise choose $R_l\to R_*$ and $h_l\downarrow0$ violating that
assertion. All roof values lie in a fixed compact interval because
$m$ is fixed. By compactness, for every neighborhood $V$ of the finite
fiber $\mathcal C_{m,\chi}(R_*)$, the corresponding vertical sections
are eventually contained in $V$. A finite union of arbitrarily short
intervals about that fiber contradicts the violation.

There is also a power modulus. Cell decomposition of the
semialgebraic neighborhood family expresses each vertical section as
a finite union of intervals with semialgebraic endpoints, after
partitioning $(R,h)$. Its length is therefore a semialgebraic function;
no assertion that arbitrary integrals of semialgebraic functions are
semialgebraic is being used. Taking its supremum over $R\in I$
preserves semialgebraicity by quantifier elimination. The resulting
bounded one-variable function tends to zero. Its graph, on a small
interval, either vanishes or is a branch of a polynomial relation.
The one-variable growth argument in the proof of
\cite[Theorem 2.12]{CosteSAG} bounds it by a positive power of $h$.
Increasing the constant treats the rest of $0<h\le1$.
\end{proof}

\begin{corollary}[A positive mass bound for the geometric remainder]
\label{cor:v62-caustic-mass}
For every fixed bounded roof interval $J$, the uninserted caustic
source satisfies
\begin{equation}\label{eq:v62-caustic-mass}
 \int_Jm^2b^{\varepsilon,\mathrm{cau}(\chi,h),1}_{n,k,m,R}(t)\,dt
 \le C(1+|J|)^{144/145}
       \bigl(D_{m,\chi}h^{\beta_{m,\chi}}\bigr)^{1/145}.
\end{equation}
Here $C$ is the constant in the inherited local weak-integrability
estimate; the new constants have the dependence just stated.
\end{corollary}
\begin{proof}
The source density is at most $P=m^2p$ after normalization and is
supported in the neighborhood in \eqref{eq:v62-caustic-tube}.
The layer-cake consequence of the inherited weak
$L^{145/144}$ bound gives, for a measurable $E\subset J$,
\[
 \int_E P\le C(1+|J|)^{144/145}|E|^{1/145}.
\]
This is Lemma~\ref{lem:v56-source-weight}. Apply it to the indicated
neighborhood intersected with $J$ and then use
Lemma~\ref{lem:v62-caustic-tube}.
\end{proof}

The last corollary is a mass estimate for an explicitly retained
physical source, not an essential-height estimate for it. Neither the
power modulus nor the separation exponent has been controlled as the
collision count grows. In particular the ordered limits required in
Corollary~\ref{cor:v51-positive-criterion} do not follow by choosing
$\varepsilon$ exponentially small in $m$: the fixed-band spectral
error has a different order of quantifiers. The present results isolate
finite-type clearance and physical caustics without replacing the
original next-collision source or assigning conditional laws to
arithmetic zero classes.
